% BHU PYQ -- Manually transcribed & error-corrected
% Paper: BCH-212 : Business Economics - I
% Exam: B.Com. (Hons.) Semester III Examination, 2016-17

\documentclass[11pt,a4paper]{article}
\usepackage[a4paper,top=2.5cm,bottom=2.5cm,left=2.8cm,right=2.8cm,headheight=14pt]{geometry}
\usepackage{amsmath,amssymb}
\usepackage[T1]{fontenc}
\usepackage[utf8]{inputenc}
\usepackage{lmodern,microtype,fancyhdr,enumitem,hyperref,lastpage,parskip}

\hypersetup{
    colorlinks=true,
    urlcolor=black,
    linkcolor=black,
    pdftitle={BCH-212: Business Economics - I -- BHU B.Com. (Hons.) Sem III 2016-17}
}

\pagestyle{fancy}
\fancyhf{}
\renewcommand{\headrulewidth}{0.4pt}
\renewcommand{\footrulewidth}{0.4pt}
\fancyhead[L]{\small   \textit{Banaras Hindu University}}
\fancyhead[R]{\small   \textit{BCH-212: Business Economics - I}}
\fancyfoot[C]{\small Page  \thepage\ of \pageref{LastPage}}


\newlist{parts}{enumerate}{2}
\setlist[parts,1]{label=(\roman*),leftmargin=*,itemsep=5pt,topsep=3pt}
\setlist[parts,2]{label=(\alph*),leftmargin=*,itemsep=3pt,topsep=2pt}

\newcommand{\pts}[1]{\hfill   \textit{[#1]}}


\begin{document}


\begin{center}
  {\large\bfseries BANARAS HINDU UNIVERSITY}\\[2pt]
  {\normalsize B.Com. (Hons.) --- Semester III Examination, 2016-17}\\[6pt]
  \rule{0.8\linewidth}{0.5pt}\\[6pt]
  {\large\bfseries Paper No. BCH-212: Business Economics - I}\\[4pt]
  {\normalsize Time: 3 Hours\hfill Maximum Marks: 70}
\end{center}

\rule{\linewidth}{0.4pt}

\smallskip



\noindent   \textit{Instructions: Answer all questions. All questions carry equal marks (14 marks each).}

\bigskip




\noindent   \textbf{Question 1.}
Discuss the characteristics of Business Economics. How does Business Economics differ from Traditional Economics?\pts{14}

\centerline{   \textbf{OR}}

Explain the following:\pts{14}

\begin{parts}
  \item Responsibilities of a Business Economist
  \item Scope of Business Economics
\end{parts}

\medskip\rule{\linewidth}{0.2pt}




\noindent   \textbf{Question 2.}
Explain the meaning and properties of an indifference curve.\pts{14}

\centerline{   \textbf{OR}}

Write a detailed note on the concept of equi-marginal utility.\pts{14}

\medskip\rule{\linewidth}{0.2pt}




\noindent   \textbf{Question 3.}
Enunciate the law of demand. Describe the factors influencing demand.\pts{14}

\centerline{   \textbf{OR}}

Define elasticity of demand. How can it be measured?\pts{14}

\medskip\rule{\linewidth}{0.2pt}

\pagebreak




\noindent   \textbf{Question 4.}
What is the importance of demand forecasting for a businessman? Discuss the steps involved in forecasting.\pts{14}

\centerline{   \textbf{OR}}

Explain the various possible approaches for forecasting the demand for an established product.\pts{14}

\medskip\rule{\linewidth}{0.2pt}




\noindent   \textbf{Question 5.}
Differentiate between the following:\pts{14}

\begin{parts}
  \item Fixed and Variable costs
  \item Historical and Replacement costs
\end{parts}

\centerline{   \textbf{OR}}

Describe clearly the external economies and diseconomies of the scale of production.\pts{14}

\vfill
\rule{\linewidth}{0.4pt}

\begin{center}\small
\href{https://bhu-exam.vercel.app/}{Click here to visit our website}\\[4pt]
\href{https://me-aryan.vercel.app/}{Click here to visit our contributor}\\[4pt]
\href{https://quashan-me.vercel.app/}{Click here to visit our contributor}
\end{center}

\end{document}
